\chapter{مقبض المغامرة}
% full version: chapters/08_nora.tex

أنت تختار أين تتغدى. أتذهب إلى المقهى المجرَّب أم تجرب مقهى جديدًا؟ إذا ذهبت دائمًا إلى المجرَّب، فلن تعرف أبدًا أن خلف الزاوية ما هو أفضل. وإذا جربت الجديد طوال الوقت، فستتغدى سيئًا كثيرًا. أين الوسط الذهبي؟ هذه المسألة، أن تبحث أو أن تنتفع بما وجدت، تواجه كل كائن يتعلم. وفي الدماغ يبدو أن المسؤولة عنها هي محطة \nt{NORA}: النورأدرينالين.

\section*{مقبض التباين}

يغيّر النورأدرينالين حدّة استجابة العصبون للإشارة. عند التأثير الضعيف يشبه العصبون مخفّت إضاءة سلسًا: إشارة أكبر قليلًا تعني استجابة أسطع قليلًا. وعند التأثير القوي يشبه مفتاحًا: تحت العتبة صمت، وفوقها استجابة كاملة. وبدقة أكبر، تبين التجارب أن النورأدرينالين يخمد الثرثرة الخلفية للعصبونات أكثر مما يخمد استجابتها للإشارة الحقيقية؛ و”مقبض التباين“ نموذج مريح لهذا الأثر.

\pencilfig{8}{\pencilart{8}}{مقبض المغامرة: أدره يسارًا فتجرب الآلة الجديد، ويمينًا فتختار المجرَّب بحزم.}

\section*{التفكير أم الحسم}

تذكّر من الفصل الثالث المجموعتين اللتين تكبح كل منهما الأخرى. حين يكون التباين منخفضًا، تعمل الاثنتان، والأضعف أخفض صوتًا قليلًا: الشبكة لا تزال ”تفكر“، تزن الخيارات. وحين يتجاوز التباين عتبة معينة، تنتقل الشبكة فجأة إلى وضع ”قررت“: تنتصر مجموعة وتصمت الأخرى. مقبض واحد ينقل الآلة بين التروي والحسم.

\section*{كم نبحث}

لنضف الضجيج. عند التباين المنخفض يرجح الضجيج بسهولة على الفرق الصغير بين الخيارات، فتجرب الآلة كثيرًا خيارًا غير الأفضل: إنها تبحث. وعند التباين العالي لا يحسم الضجيج شيئًا، فتأخذ الآلة دائمًا ما تراه الأفضل: إنها تنتفع.

في نموذجنا يتغير العالم من حين لآخر: يصير الخيار الأفضل أسوأ. والربح يتبع المقبض على شكل حرف U مقلوب. تباين أقل مما ينبغي: الآلة تهيم على غير هدى. تباين أكثر مما ينبغي: لا تلاحظ أن العالم تغير، وتذهب بعناد إلى المقهى الذي فسد. وكلما تغير العالم أكثر، انخفض الضبط الأفضل: في عالم متقلب يربح من يبحث.

\pencilfig{108}{%
% points: models/nora_explore.py, reward as a share of the best possible, gain 0.5 ... 64 (doubling); the y axis starts at 0.70, hence the break mark
\draw[pen] (0,0) -- (6.3,0); \draw[pen] (0,0) -- (0,0.18); \draw[pen] (0,0.34) -- (0,2.4);
\draw[pcGray, line width=0.6pt] (-0.1,0.14) -- (0.1,0.22); \draw[pcGray, line width=0.6pt] (-0.1,0.3) -- (0.1,0.38);
\node[lbl, anchor=north] at (3.0,-0.05) {مقبض التباين};
\node[lbl, anchor=south, rotate=90] at (-0.1,1.3) {الربح};
\draw[penlite=pcBlue] plot[smooth] coordinates {(0.00,0.55) (0.85,0.79) (1.70,1.19) (2.55,1.68) (3.40,2.00) (4.25,2.07) (5.10,2.01) (5.95,1.91)};
\draw[penlite=pcRed] plot[smooth] coordinates {(0.00,0.55) (0.85,0.69) (1.70,0.93) (2.55,1.24) (3.40,1.40) (4.25,1.28) (5.10,1.06) (5.95,0.89)};
\draw[dotpen=pcBlue, wash=pcBlue] (4.25,2.07) circle (0.07);
\draw[dotpen=pcRed, wash=pcRed] (3.40,1.40) circle (0.07);
\node[lbl, text=pcBlue, anchor=south] at (4.25,2.15) {العالم يتغير نادرًا};
\node[lbl, text=pcRed, anchor=north] at (3.40,1.30) {كثيرًا};
\node[lbl, anchor=north west] at (0.05,-0.35) {تهيم بلا هدى}; \node[lbl, anchor=north east] at (6.3,-0.35) {لا تلاحظ التغيير};
}{ربح النموذج عند مواضع مختلفة للمقبض. لكل منحنى قمة؛ وحين يتغير العالم كثيرًا تكون أدنى وإلى اليسار.}

\section*{من يدير المقبض}

العلامة الطبيعية على التغيير هي المفاجأة: صارت الأخطاء أكبر من المعتاد. ووفق إحدى الفرضيات، عن هذا التغيير المفاجئ بالذات يخبر النورأدرينالين. وهنا دقيقة مهمة: المستوى الخلفي الثابت للنورأدرينالين يرتبط أكثر بالبحث والتشتت، أما الرشقات القصيرة عند حدث مهم فترتبط بالتركيز والحسم. ربما تعمل الرشقة كالمقاطعة في الحاسوب: ”اترك ما بيدك، العالم تغير، أعد بناء الخطة“.

وعلامة غير مباشرة على هذا النظام تُرى من الخارج: تتسع الحدقة حين يتغير شيء ما على غير توقع ويبدأ الناس في التعلم أسرع. لكن الحدقة لا تتبع النورأدرينالين وحده، فهي تلميح لا جهاز قياس.

واعتراف صريح: في نموذجنا لم يربح ضبط المقبض بحسب المفاجأة إلا قليلًا جدًا مقارنة بأفضل ضبط ثابت. أما أين يؤتي هذا الضبط ثماره تحديدًا فلا يزال علينا أن نفهمه.

\fullversion{8}
